\documentclass[
	11pt,oneside,a4paper,
	headsepline,footsepline,
	%plainfootsepline,plainheadsepline,
	fleqn,
	%flushbottom,
	%raggedbottom
]{memoir}
%]{scrartcl}
%]{article}
%\usepackage[a4paper, total={6in, 8in}]{geometry}

%------------------------------------------
% FLAG DE BORRADOR
%------------------------------------------
%\def\borrador{}
%------------------------------------------

%------------------------------------------
% CONFIGURACION DEL DOCUMENTO
%------------------------------------------

\setlength {\marginparwidth }{2cm}

%------------------------------------------
% INCLUSION DE PAQUETES COMUNES
%------------------------------------------
%\usepackage[]{fontspec}
%\setmainfont{Times New Roman}
\usepackage[]{newtx}
\usepackage[spanish]{babel}

% ======================================================
% GESTION DE COLORES Y GRAFICOS
% ======================================================
\usepackage[dvipsnames,svgnames]{xcolor}
\usepackage{graphicx}
% ======================================================
% MEJORA LA LEGIBILIDAD GENERAL DEL DOCUMENTO Y LAYOUT
% ======================================================
% Microtype tiene que preceder a Ragged2e
\usepackage[babel=true,expansion=true,protrusion=true]{microtype}
% Ragged2e tiene que ir despues de microtype para funcionar correctamente sino cambia las alineaciones y no quedan bien
\usepackage{ragged2e}

% ======================================================
% MARCAS DE AGUA
% ======================================================
\usepackage{draftwatermark}
% ======================================================
% PERMITE INCORPORAR HIPERVINCULOS (internos y URL externas)
% ======================================================
\usepackage[
	colorlinks,
	urlcolor = blue,
	linkcolor = black,
	citecolor = blue
	%pdfborder = {0 0 0}
]{hyperref}


% ======================================================
% MEJORES LISTAS DE ENUMERACION
% ======================================================
\usepackage{enumitem}
% Para poder tachar texto con \uwave{}, uwave{}, sout{}
\usepackage[normalem]{ulem} % El parametro normalem es para no alterar el emphasis (italica)

% ======================================================
% TODOLIST
% ======================================================
\setlength{\marginparwidth}{2cm}
\usepackage[colorinlistoftodos,textsize=tiny]{todonotes}
\ifdefined\borrador
	\newcommand{\revisar}[1]{\todo[backgroundcolor=yellow!50]{\textbf{Revisar:} #1}}
	\newcommand{\ampliar}{\todo[backgroundcolor=blue!30]{Ampliar}}
	\newcommand{\revisarcita}{\colorbox{red}{\textcolor{white}{\ CITA\ }}}
	\newcommand{\sincita}{\colorbox{red}{\textcolor{white}{\ CITA\ }}}
	\newcommand{\faltafigura}{\missingfigure{Falta figura\ldots}}
	\newcommand{\citaurl}{\todo[backgroundcolor=green!50]{URL}}
\else
	\newcommand{\revisar}[1]{}
	\newcommand{\ampliar}{}
	\newcommand{\revisarcita}{}
	\newcommand{\sincita}{}
	\newcommand{\faltafigura}{}
	\newcommand{\citaurl}{}
\fi

\usepackage{csquotes}

% ======================================================
% GANTT
% ======================================================
\usepackage{pgfgantt}
% ======================================================
% FLOATS y REORDENAMIENTO DE TEXTO ALREDEDOR DE FIGURAS
% ======================================================
\usepackage{float}
\usepackage{wrapfig}
\usepackage{csquotes}
\usepackage[]{enumitem}
%------------------------------------------
% CONFIGURACION DE BIBLIOGRAFIA
%------------------------------------------
\usepackage[
	backend=biber,
	natbib=true,
	sorting=anyt,  %% Alfabetico
	%sorting=none,	%% En orden de aparición
	%style=nature,
	style=authoryear-ibid,
	%style=numeric,
	%style=unsrt
	%citestyle=numeric
]{biblatex}
\addbibresource{refs_used.bib}
%\nocite{*}

%\DeclareBibliographyCategory{cited}
%\AtEveryCitekey{\addtocategory{cited}{\thefield{entrykey}}}
%\addbibresource{refs_used.bib}
%\BiblatexSplitbibDefernumbersWarningOff

\urlstyle{sf}

\ifdefined\borrador
	\SetWatermarkText{Borrador}
	\SetWatermarkScale{0.8}
	\SetWatermarkLightness{0.9}
\else
	\SetWatermarkText{}
\fi

%------------------------------------------
% FUENTES
%------------------------------------------
%\usepackage[sfdefault,scaled=1]{FiraSans}
%\usepackage[libertine,scaled=1]{newtx}
\setmainfont[Mapping=tex-text]{Linux Libertine O}
%------------------------------------------

\graphicspath{
	{figures/}
	{images}
}

%\clearpairofpagestyles
%\ihead{\today}								% Header interno: Fecha
% ENCABEZADO
%\ohead{\textit{Plan de Tesis}}		% Header interno: Titulo
%\ihead{\raisebox{-7mm}{\includegraphics[height=10mm]{./images/utn-black}}}		% Footer interno: 
% PIE DE PAGINA
%\ofoot{Página \pagemark}	% Header externo: Evento
%Titulo
%\ohead{\raisebox{-10mm}{\includegraphics[height=10mm]{./images/MISI3}}}

% INTERLINEADO 1 1/2 
%\OnehalfSpacing

\pgfplotsset{compat=1.18}
\renewcommand\fbox{\fcolorbox{gray!50}{white}}

\usepackage{pdflscape}

\begin{document}

\thispagestyle{empty}

\begin{center}
	\vspace*{3cm}
	
	\HUGE{Caracterización estructural de formaciones astronómicas, utilizando un enfoque de grafos, para reconocimiento de cumulos estelares en galaxias cercanas}
	
	\vspace{2cm}
	
	\LARGE{Universidad Tecnológica Nacional\\
		Facultad Regional Córdoba}
	
	\vspace{2cm}
	
	\LARGE{Proyecto de Investigación y Desarrollo (PID)}
	
	\vspace{2cm}
\end{center}

\begin{flushright} 
	\Large{
		Tesista: Esp. Ing. Martin Casatti\\
		Director: Dr. Oscar Medina\\
		CoDirectors: Mgr. Cynthia Corso
	}
\end{flushright} 


\newpage

\ifdefined\borrador
	\listoftodos
\fi 

\section*{Plan de cursos}

Se detalla en la siguiente sección el plan de cursos a realizar y cursos ya realizados (aquellos que tienen nota y fecha de aprobación).

El tipo de curso se define según la siguientes siglas:
\begin{description}
    \item[PF] PROFORVIN
    \item[PG] POSGRADO
    \item[MI] MAESTRÍA EN INGENIERÍA  
\end{description}

\begin{table}[H]
    \centering
    \begin{tabular}{|l|c|l|c|c|}
    \hline
    \textbf{Nombre} &
      \textbf{Carga horaria} &
      \textbf{Aprobado} &
      \textbf{Tipo} &
      \textbf{Nota} \\ \hline
    \begin{tabular}[c]{@{}l@{}}Introducción a la Investigación, \\ el Desarrollo y la Innovación.\end{tabular}                & 64 & 29/06/2017 & PF & 9                     \\ \hline
    \begin{tabular}[c]{@{}l@{}}Redes Neuronales y Lógica Difusa \\ en Ingeniería\end{tabular} &
      60 &
      05/03/2021 &
      PG &
      10 \\ \hline
    Introducción a la Ciencia de Datos &
      60 &
      02/10/2020 &
      PG &
      10 \\ \hline
    \begin{tabular}[c]{@{}l@{}}Métodos Empíricos en Ingeniería \\ de Software\end{tabular} &
      60 &
      2018/2019 &
      PG &
       \\ \hline
    \begin{tabular}[c]{@{}l@{}}Modelos de Organizaciones y Sistemas \\ de Información\end{tabular} &
      60 &
      19/05/2016 &
      MI &
      8 \\ \hline
    Ingeniería de Software(Elect.) &
      60 &
      29/07/2016 &
      MI &
      9 \\ \hline
    \begin{tabular}[c]{@{}l@{}}Modelado Conceptual de Sistemas \\ de Información(Elec.)\end{tabular} &
      60 &
      31/10/2016 &
      MI &
      9 \\ \hline
    \begin{tabular}[c]{@{}l@{}}C.A.P. Estimaciones de Software de \\ Sistemas de Información\end{tabular} &
      60 &
      14/11/2016 &
      MI &
      10 \\ \hline
    Seminario Integrador &
       &
      18/05/2018 &
      MI &
      10 \\ \hline
    Trabajo Final Integrador &
       &
      18/05/2018 &
      MI &
      10 \\ \hline
    \begin{tabular}[c]{@{}l@{}}Tópicos de Machine Learning aplicados\\ a Big Data\end{tabular} &
      36 &
       &
      PG &
      \multicolumn{1}{l|}{} \\ \hline
    \begin{tabular}[c]{@{}l@{}}Aspectos teóricos y metodológicos del \\ análisis computacional de redes sociales\end{tabular} & 60 &            & PG & \multicolumn{1}{l|}{} \\ \hline
    \end{tabular}
    \end{table}


\end{document}